% Compile with pdflatex; the cropped PDF is ready for \includegraphics.
% These are SCHEMATIC GLOBAL views, not coordinate plots of the quartics.
% The connected contours display the correct local contact orders and
% ordinary nodes. The center triangles show schematic incidence data.
% Drawing functions below are chosen only to obtain these local types:
% y = +/- f(x) sqrt(1-(x/2.7)^2). They are NOT equations of C_i.
% For general smooth auxiliary centers, both transformed curves have degree 4.
% Contracting A_i B_i creates one node from two residual transverse points;
% contracting P_i A_i or P_i B_i creates a smooth point. Choose the second
% triangle to avoid N1, so N1 persists as N1'.
\documentclass[tikz,border=5pt]{standalone}
\usepackage{amsmath,amssymb}
\usepackage[T1]{fontenc}
\usepackage{lmodern}
\usetikzlibrary{arrows.meta}
\definecolor{ink}{HTML}{193347}
\definecolor{branchA}{HTML}{156E9A}
\definecolor{branchB}{HTML}{C06435}
\definecolor{muted}{HTML}{667681}
\definecolor{pale}{HTML}{F1F5F7}
\pgfmathdeclarefunction{fzero}{1}{\pgfmathparse{2.3704*(#1)^3/(.55+(#1)^2)}}
\pgfmathdeclarefunction{fone}{1}{\pgfmathparse{3.1316*((#1)+.65)^2*((#1)-1.35)/(1+(#1)^2)}}
\pgfmathdeclarefunction{ftwo}{1}{\pgfmathparse{3.5577*(#1)*((#1)^2-1.69)/(.65+(#1)^2)}}
\begin{document}
\begin{tikzpicture}[x=1cm,y=1cm,>=Latex,text=ink,
        every node/.style={font=\small},
        a/.style={branchA,line width=1.25pt},
        b/.style={branchB,line width=1.25pt},
        panel/.style={rounded corners=4pt,fill=pale,draw=none},
        triangle/.style={muted!85,dashed,line width=.8pt},
        pointlabel/.style={fill=pale,inner sep=1.5pt}]
    \node[font=\large\bfseries] at (8.2,5.85) {Two quadratic transformations};
    \node at (8.2,5.3) {$C_0:\ (YZ-X^2)^2-Y^3Z=0$};
    \node[font=\footnotesize,muted] at (8.2,4.8)
    {Schematic global views: all singularities shown on each connected curve};
    \foreach \x in {0,8.2,16.4}{
            \path[panel] (\x-3.35,-4.65) rectangle (\x+3.35,4.25);
        }
    % Stage 0: third-order contact at the single singular point.
    \begin{scope}
        \node[font=\normalsize\bfseries] at (0,3.85) {$C_0$: triple contact $A_5$};
        \coordinate (P0) at (0,0);
        \coordinate (A0) at (-1.9,{-fzero(-1.9)*sqrt(1-1.9^2/2.7^2)});
        \coordinate (B0) at (1.65,{fzero(1.65)*sqrt(1-1.65^2/2.7^2)});
        \draw[triangle] (P0)--(A0)--(B0)--cycle;
        \draw[a,domain=0:180,samples=361,variable=\t]
        plot ({2.7*cos(\t)},{fzero(2.7*cos(\t))*sin(\t)});
        \draw[b,domain=180:360,samples=361,variable=\t]
        plot ({2.7*cos(\t)},{fzero(2.7*cos(\t))*sin(\t)});
        \foreach \p in {P0,A0,B0}{\fill (\p) circle (2pt);}
        \node[pointlabel,below=6pt] at (P0) {$P_0$};
        \node[pointlabel,above left=5pt] at (A0) {$A_0$};
        \node[pointlabel,above=10pt] at (B0) {$B_0$};
        \node[font=\footnotesize] at (0,-3.55) {$\varphi_1\text{ centered at }\{P_0,A_0,B_0\}$};
        \node at (0,-4.17) {$\operatorname{Sing}(C_0)=\{P_0\}$};
    \end{scope}
    \draw[->,line width=1pt] (3.5,0)--(4.7,0);
    \node at (4.1,.5) {$\varphi_1$};
    % Stage 1: second-order tangency and a transverse ordinary node.
    \begin{scope}[shift={(8.2,0)}]
        \node[font=\normalsize\bfseries] at (0,3.85) {$C_1$: tacnode $+$ one node};
        \coordinate (P1) at (-.65,0);
        \coordinate (N1) at (1.35,0);
        \coordinate (A1) at (-1.95,{-fone(-1.95)*sqrt(1-1.95^2/2.7^2)});
        \coordinate (B1) at (.35,{-fone(.35)*sqrt(1-.35^2/2.7^2)});
        \draw[triangle] (P1)--(A1)--(B1)--cycle;
        \draw[a,domain=0:180,samples=401,variable=\t]
        plot ({2.7*cos(\t)},{fone(2.7*cos(\t))*sin(\t)});
        \draw[b,domain=180:360,samples=401,variable=\t]
        plot ({2.7*cos(\t)},{fone(2.7*cos(\t))*sin(\t)});
        \foreach \p in {P1,N1,A1,B1}{\fill (\p) circle (2pt);}
        \node[pointlabel,below=7pt] at (P1) {$P_1$};
        \node[pointlabel,right=8pt] at (N1) {$N_1$};
        \node[pointlabel,above left=5pt] at (A1) {$A_1$};
        \node[pointlabel,above=10pt] at (B1) {$B_1$};
        \node[font=\footnotesize] at (0,-3.22) {$N_1=q_1(\overline{A_0B_0})$};
        \node[font=\footnotesize] at (0,-3.65) {$q_2\text{ centered at }\{P_1,A_1,B_1\}$};
        \node at (0,-4.17) {$\operatorname{Sing}(C_1)=\{P_1,N_1\}$};
    \end{scope}
    \draw[->,line width=1pt] (11.7,0)--(12.9,0);
    \node at (12.3,.5) {$\varphi_2$};
    % Stage 2: three ordinary nodes on one connected curve, without insets.
    \begin{scope}[shift={(16.4,0)}]
        \node[font=\normalsize\bfseries] at (0,3.85) {$C_2$: three ordinary nodes};
        \draw[a,domain=0:180,samples=451,variable=\t]
        plot ({2.7*cos(\t)},{ftwo(2.7*cos(\t))*sin(\t)});
        \draw[b,domain=180:360,samples=451,variable=\t]
        plot ({2.7*cos(\t)},{ftwo(2.7*cos(\t))*sin(\t)});
        \coordinate (N1p) at (-1.3,0);
        \coordinate (P2) at (0,0);
        \coordinate (N2) at (1.3,0);
        \foreach \p in {N1p,P2,N2}{\fill (\p) circle (2pt);}
        \node[pointlabel,right=8pt] at (N1p) {$N_1'$};
        \node[pointlabel,right=8pt] at (P2) {$P_2$};
        \node[pointlabel,right=8pt] at (N2) {$N_2$};
        \node[font=\footnotesize] at (0,-3.22) {$N_1'=\varphi_2(N_1)$};
        \node[font=\footnotesize] at (0,-3.65) {$N_2=\varphi_2(\overline{A_1B_1})$};
        \node at (0,-4.17) {$\operatorname{Sing}(C_2)=\{P_2,N_1',N_2\}$};
    \end{scope}
\end{tikzpicture}
\end{document}